\section{2026 年 11 月核验增补}

本节按论文首次提交时间排列。每项方法附有可核验的外部入选依据。

% METHOD: 2609.07174
\paragraph{PhysMAS：按部件组织材料与运动再做共享域仿真}
2026.09~\cite{2609_07174}。PhysMAS 接收运动描述、四张已标定图像和已对齐的高斯场景。场景模块保持物体及部件身份，材料模块为各部件给出物理配置，再用共享域 MPM 仿真生成动态高斯，避免为每个场景反复做生成梯度优化。

论文比较六类材料场景的语义对齐、感知的物理合理性及参考视频运动一致性。可选本构配置和空间分辨率限制了部件行为与长时互动的覆盖，遮挡下的部件关联也仍可能不稳定。

\textbf{入选依据} 论文署名单位包括 Meituan（\href{https://arxiv.org/abs/2609.07174}{论文原文}）。

% METHOD: 2610.13945
\paragraph{Splat2Act：从实景高斯生成面向部署环境的 VLA 示范}
2026.10~\cite{2610_13945}。Splat2Act 用手机拍摄室内场景并另采目标物体，重建可渲染环境并自动识别候选容器。任务生成和示范生成器在仿真中构造抓取数据，用来将现成的视觉语言动作策略适配给定地点、容器和物品。

作者在办公室及会议室测试训练中没有出现的目标位置组合，测试任务仍使用已知物体身份，并根据示范生成器确认可行状态。结果针对所测部署环境和取物任务。作者页可核实 Waymo LLC 和 Toyota Research Institute 的署名。

\textbf{入选依据} 论文署名单位包括 Waymo、Toyota Research Institute（\href{https://arxiv.org/html/2610.13945v1}{论文原文}）。

% METHOD: 2611.20283
\paragraph{CiGS：从高斯重建表面求解固体内部温度场}
2026.11~\cite{2611_20283}。CiGS 用已配准的可见光和热图重建高斯表面，再提取带热边界观测的点云，直接求解边界积分方程以推断内部温度。深度斜率过滤剔除跨深度断裂的采样点，内部积分对密度常数项单独计算，整个过程不需要表面或体积网格。

作者比较完整覆盖和单视图采样下的精度及耗时，并使用无矩阵的可微求解。由边界观测推断内部温度依赖重建表面和正确的物理边界设定，单视图结果尤其受表面覆盖限制。

\textbf{入选依据} 论文署名单位包括 Google DeepMind（\href{https://arxiv.org/abs/2611.20283}{论文原文}）。
